\subsection{An Alternate Method: Reverse Multiplication}
This alternate method reverses the area method for multiplication.
\begin{Example}We begin with an example.
\[
(2x^5-9x^4+19x^3-20x^2+7x+3)\div (2x^2-5x+1)
\]
\begin{lpic}{}{5}
\end{lpic}
\end{Example}
\begin{enumerate}
\item First, guess the degree of the result. In this example, The degree of the dividend is \makeblank[.5in] and the degree of the divisor is \makeblank[.5in]. The degree of the quotient will be \makeblank.
\item The number of columns we need is two more than the degree of the quotient. In this case we need \makeblank columns.
\item We need a row for each term of the divisor (including zeros). In this example we need \makeblank[.5in] rows.
\item Label the columns with descending degrees. The second-last column is the constant. The last column is the remainder.
\item Write the divisor down the left, in descending order.
\item Write the first term of the dividend in the upper left corner.
\item\label{pdiv_startloop} Reverse multiplication to figure out what term goes above it.
\item Multiply to fill in the column.
\item\label{pdiv_endloop} Each diagonal must add up to the corresponding term. Figure out what number needs to go in the next empty box in the top row to get the correct sum.
\item Repeat steps \ref{pdiv_startloop} -- \ref{pdiv_endloop} until you have filled in the top number in the remainder column.
\item Fill in the rest of the remainder column to get the appropriate totals.
\end{enumerate}